%%-*-latex-*-

\documentclass[11pt]{article}

\usepackage{pst-tree,pst-eps}
\usepackage{amsmath}

\newcommand\fun[1]{\textsf{#1}}

\begin{document}

\TeXtoEPS
\pstree[nodesep=2pt,levelsep=17pt,treesep=14pt]{\TR{\fun{cat}}}{
  \TR{\texttt{[]}}
  \TR{\rnode[r]{Qa}{\(t\)}}
}
\quad\raisebox{-3pt}{\(\xrightarrow{\alpha}\)}\quad
\pstree[levelsep=17pt,treesep=14pt]{\TR{\rnode[cb]{Qup}{\(\circ\)}}}{}
\qquad\qquad
\pstree[nodesep=2pt,levelsep=17pt,treesep=14pt]{\TR{\fun{cat}}}{
  \pstree{\TR{\texttt{|}}}{
    \TR{\rnode[br]{I}{\(x\)}}
    \TR{\rnode[br]{P}{\(s\)}}
  }
  \TR{\rnode[br]{Qb}{\(t\)}}
}
\quad\raisebox{-3pt}{\(\xrightarrow{\beta}\)}\quad
\pstree[nodesep=2pt,levelsep=15pt,treesep=14pt]%
  {\TR{\rnode[bl]{cons}{\texttt{|}}}}{
  \Tn
  \Tn
  \pstree{\TR{\rnode[bl]{joinP}{\rnode[br]{joinQ}{\TR{\fun{cat}}}}}}{
    \Tn
  }
}
\nccurve[angleA=-90,angleB=0,nodesepB=1pt,nodesepA=-1pt]{->}{Qup}{Qa}
\nccurve[angleA=-130,angleB=-55,nodesepA=2pt,nodesepB=1pt,ncurvB=1]{->}{cons}{I}
\nccurve[angleA=-120,angleB=-45,nodesepA=5pt]{->}{joinP}{P}
\nccurve[angleA=-60,angleB=-50,nodesepA=5pt,nodesepB=1pt,ncurvA=1]{->}{joinQ}{Qb}
\endTeXtoEPS

\end{document}
